\chapter{Exponential Sums over Primes}
\begin{center}\itshape
Multiply the primes by an irrational number and keep the fractional parts:\\
they fill the unit interval evenly. Vinogradov proved it in 1937,\\
and the same sum, $\sum\Lambda(n)e(n\alpha)$, carries the three primes theorem.
\end{center}

Weyl's theorem of 1916 (Tome~5, Theorem~11.4) says that $n\alpha$, $n^2\alpha$, and every polynomial with an irrational coefficient, are equidistributed modulo $1$. The proof differences the polynomial until it becomes linear. The primes cannot be differenced: $p+h$ is rarely prime. For twenty years nobody could show that $p\alpha$ is equidistributed, although it is ``obviously'' true. Ivan Matveevich Vinogradov did it in 1937, with a combinatorial decomposition of the primes and a bound for bilinear sums, and in the same paper proved that every large odd number is a sum of three primes (Chapter~7). Robert Vaughan's identity of 1977, which we used in Chapter~4, turns Vinogradov's argument into three pages.

Exponential sums with an algebraic phase behave differently. In 1926 Hendrik Kloosterman, studying the representations of integers by quadratic forms in four variables, met the sums $\sum_xe(\frac{ax+b\overline x}p)$, where $\overline x$ is the inverse of $x$ modulo $p$, and showed that they are at most a constant times $p^{3/4}$, not $p$. André Weil proved the bound $2\sqrt p$ in 1948, as a consequence of the Riemann hypothesis for curves over finite fields; it is far beyond this book, but its cousin for Salié sums is a pleasant exercise in Gauss sums, and we prove it exactly.

\begin{reminder}[from Tomes 4 and 5, and Chapters 2 and 4]
\emph{Equidistribution} (Tome~4, Theorem~6.13). A sequence $(u_n)$ is equidistributed modulo $1$ if and only if $\frac1N\sum_{n\le N}e(ju_n)\to0$ for every integer $j\ne0$; here $e(t)=e^{2\pi it}$ and $\|t\|$ is the distance from $t$ to the nearest integer.
\emph{Geometric sums} (Tome~5, Lemma~11.1). $|\sum_{M<n\le M+N}e(\beta n)|\le\min(N,\frac1{2\|\beta\|})$.
\emph{Sums of minima} (Tome~5, Lemma~11.6). If $|\alpha-\frac aq|\le\frac1{q^2}$ with $(a,q)=1$, then for $X,N\ge1$,
$\sum_{1\le x\le X}\min(N,\frac1{\|\alpha x\|})\le(\frac Xq+1)(4N+2q+2q\log q)$.
\emph{Gauss sums} (Tome~5, Theorem~3.5 and Exercise~3.3). For the Legendre symbol modulo an odd prime $p$, $\sum_x(\frac xp)e(\frac{cx}p)=(\frac cp)g_p$ for every $c$, where $g_p=\sum_x(\frac xp)e(\frac xp)$ satisfies $g_p^2=(\frac{-1}p)p$.
\emph{Primes}. Vaughan's identity (Lemma~\ref{lem:vaughan}), the bounds $\sum_{m\le N}\tau(m)^2\le N(1+\log N)^4$ (Lemma~\ref{lem:phisum}), and Siegel--Walfisz (Corollary~\ref{cor:sw}): $|\psi(x;q,a)-\frac x{\varphi(q)}|\le C(A)\frac x{(\log x)^A}$ for $q\le(\log x)^A$, $(a,q)=1$.
\end{reminder}

\section{Rational approximations}

Every estimate in this chapter depends on how well $\alpha$ is approximated by fractions. We start with Dirichlet's pigeonhole and a variant of the lemma on sums of minima.

\begin{lemma}[Dirichlet]\label{lem:dirichlet}
For real $\alpha$ and $Q\ge1$ there are coprime integers $a$ and $q$ with $1\le q\le Q$ and $|\alpha-\frac aq|\le\frac1{qQ}$.
\end{lemma}
\begin{proof}
Let $n=\lfloor Q\rfloor$. Cut $[0,1]$ into $n+1$ closed intervals of length $\frac1{n+1}$. The $n+2$ numbers $\{j\alpha\}$, $0\le j\le n$, and $1$ lie in them, so two lie in the same one; their difference gives integers $1\le q\le n$ and $a$ with $|q\alpha-a|\le\frac1{n+1}<\frac1Q$. Dividing $a$ and $q$ by their gcd only decreases $q$, and keeps $|\alpha-\frac aq|\le\frac1{qQ}$.
\end{proof}

\begin{lemma}\label{lem:summin2}
Let $|\alpha-\frac aq|\le\frac1{q^2}$ with $(a,q)=1$, $q\ge1$, and let $M,N\ge1$. Then
\[
\sum_{m\le M}\min\Bigl(\frac Nm,\frac1{\|\alpha m\|}\Bigr)\le40\Bigl(\frac Nq+M+q\Bigr)\log(3qM).
\]
\end{lemma}
\begin{proof}
Write $m=jq+r$ with $0\le r<q$, and fix $j$. With $\beta=\alpha jq$, $\alpha m=\beta+\frac{ar}q+\delta_r$, where $|\delta_r|=r|\alpha-\frac aq|<\frac1q$. As $r$ runs over $0,\dots,q-1$, the points $\beta+\frac{ar}q$ run over the $q$ points $\beta+\frac sq$ modulo $1$. If $\frac kq\le\|\alpha m\|<\frac{k+1}q$, then $\beta+\frac{ar}q$ lies at distance less than $\frac1q$ from the two arcs where $\|\cdot\|\in[\frac kq,\frac{k+1}q)$, that is in two arcs of length $\frac3q$, which contain at most $8$ of the points. So for each $k\ge0$ at most $8$ values of $r$ have $\|\alpha m\|\in[\frac kq,\frac{k+1}q)$.
For $j\ge1$, every $m$ in the block has $\frac Nm\le\frac N{jq}$; bounding the terms with $k=0$ by $\frac N{jq}$ and the others by $\frac qk$, the block contributes at most $\frac{8N}{jq}+8q\sum_{k\le q/2}\frac1k\le\frac{8N}{jq}+8q(1+\log q)$.
For $j=0$ and $1\le m\le\frac q2$, $\|\frac{am}q\|\ge\frac1q$ and $|\alpha m-\frac{am}q|\le\frac m{q^2}\le\frac1{2q}$, so $\|\alpha m\|\ge\frac12\|\frac{am}q\|$; the residues $am$ are distinct and non-zero, so these $m$ contribute at most $2\sum_{m\le q/2}\|\frac{am}q\|^{-1}\le4q(1+\log q)$. The $m\in\left]\frac q2,q\right[$ have $\frac Nm<\frac{2N}q$ and contribute at most $\frac{16N}q+8q(1+\log q)$ by the counting above.
The blocks $j\ge1$ that meet $[1,M]$ have $j\le\frac Mq$. Adding up, with $\sum_{j\le M/q}\frac1j\le1+\log M$, the sum is at most
\[
\frac{16N}q+12q(1+\log q)+\frac{8N}q(1+\log M)+8M(1+\log q),
\]
which is less than $40(\frac Nq+M+q)\log(3qM)$, since $1+\log q$ and $1+\log M$ are at most $\log(3qM)$.
\end{proof}

\section{Vinogradov's estimate}

Let $S(\alpha)=\sum_{n\le N}\Lambda(n)e(n\alpha)$. At $\alpha=0$ it is $\psi(N)\sim N$; at $\alpha=\frac12$ it is $-\psi(N)$ up to the powers of $2$. The theorem says that $S(\alpha)$ is small whenever $\alpha$ is close to a fraction with a large denominator.

\begin{theorem}[Vinogradov, 1937]\label{thm:vinogradov3}
Let $N\ge3$, and $|\alpha-\frac aq|\le\frac1{q^2}$ with $(a,q)=1$. Then
\[
\sum_{n\le N}\Lambda(n)e(n\alpha)\ll\Bigl(\frac N{\sqrt q}+N^{4/5}+\sqrt{Nq}\Bigr)(\log N)^5,
\]
with an absolute implied constant.
\end{theorem}
\begin{proof}
If $q>N$ the right side exceeds $N$ and there is nothing to prove; so let $q\le N$, and write $\mathcal L=\log N$. Take $U=V=N^{2/5}$ in Lemma~\ref{lem:vaughan}, and split $S(\alpha)=S_1+S_2+S_3+S_4$ accordingly.

\emph{Type I.} $|S_1|\le\psi(U)\ll U$. Next,
\[
S_2=\sum_{b\le V}\mu(b)\sum_{d\le N/b}(\log d)e(\alpha bd),
\]
and Abel summation with the bound for geometric sums gives $|\sum_{d\le D}(\log d)e(\beta d)|\le2\log D\cdot\min(D,\frac1{2\|\beta\|})$. So
\[
|S_2|\le2\mathcal L\sum_{b\le V}\min\Bigl(\frac Nb,\frac1{\|\alpha b\|}\Bigr).
\]
In $S_3=-\sum_{t\le UV}\lambda_t\sum_{d\le N/t}e(\alpha td)$, $|\lambda_t|\le\log t$, so $|S_3|\le\mathcal L\sum_{t\le UV}\min(\frac Nt,\frac1{\|\alpha t\|})$. By Lemma~\ref{lem:summin2} with $M=UV=N^{4/5}$, and $\log(3qN^{4/5})\le3\mathcal L$,
\[
|S_1|+|S_2|+|S_3|\ll\Bigl(\frac Nq+N^{4/5}+q\Bigr)\mathcal L^2.
\]
\emph{Type II.} $S_4=\sum_{km\le N,\,k>U,\,m>V}\Lambda(k)c(m)e(\alpha km)$. Cut the range of $m$ into at most $2\mathcal L$ blocks $\left]M,2M\right]$ with $V\le M<\frac NU$. On one block, by Cauchy--Schwarz,
\[
\Bigl|\sum_{M<m\le2M}c(m)\sum_{U<k\le N/m}\Lambda(k)e(\alpha km)\Bigr|^2\le\sum_{m\le2M}\tau(m)^2\cdot\Sigma,
\]
where $\Sigma=\sum_{M<m\le2M}|\sum_{U<k\le N/m}\Lambda(k)e(\alpha km)|^2$. Expanding the square and summing over $m$ first, for each pair $k,k'\le\frac NM$ the $m$ run over an interval of length at most $M$, so by the bound for geometric sums and $\Lambda(k)\Lambda(k')\le\mathcal L^2$,
\[
\Sigma\le\mathcal L^2\sum_{k,k'\le N/M}\min\Bigl(M,\frac1{2\|\alpha(k-k')\|}\Bigr)\le\mathcal L^2\frac NM\Bigl(M+2\sum_{1\le h\le N/M}\min\Bigl(M,\frac1{\|\alpha h\|}\Bigr)\Bigr).
\]
By Tome~5, Lemma~11.6 with $X=\frac NM$, the last sum is at most $(\frac N{Mq}+1)(4M+2q+2q\log q)\le(\frac{4N}q+4M+\frac{2N}M+2q)(1+\log q)$. With $\sum_{m\le2M}\tau(m)^2\le2M(1+\log2M)^4$ and $\log q\le\mathcal L$, the square of the block is therefore
\[
\ll\mathcal L^7\Bigl(\frac{N^2}q+NM+\frac{N^2}M+Nq\Bigr).
\]
Taking square roots, and using $M<\frac NU$ for the second term and $M\ge V$ for the third, the block is $\ll\mathcal L^{7/2}(\frac N{\sqrt q}+\frac N{\sqrt U}+\frac N{\sqrt V}+\sqrt{Nq})$, and there are at most $2\mathcal L$ blocks. As $\frac N{\sqrt U}=\frac N{\sqrt V}=N^{4/5}$, and $\frac Nq\le\frac N{\sqrt q}$, $q\le\sqrt{Nq}$, the theorem follows.
\end{proof}

When $q$ is about $\sqrt N$, the bound is $N^{4/5}(\log N)^5$, a saving of a power of $N$ over the trivial bound $\psi(N)$. When $q$ is small, the theorem says nothing, and nothing like it can be true: near $\frac aq$, $S(\alpha)$ is large. The next lemma computes it there, from the prime number theorem in progressions.

\begin{lemma}[major arcs]\label{lem:majorarc}
For all $A,B>0$ there is a constant $C(A,B)$ such that, for $N\ge3$, $1\le q\le(\log N)^B$, $(a,q)=1$ and real $\beta$,
\[
\Bigl|\sum_{n\le N}\Lambda(n)e\Bigl(n\Bigl(\frac aq+\beta\Bigr)\Bigr)-\frac{\mu(q)}{\varphi(q)}\sum_{n\le N}e(\beta n)\Bigr|\le C(A,B)(1+|\beta|N)\frac N{(\log N)^A}.
\]
The constant cannot be computed by the proof.
\end{lemma}
\begin{proof}
First, the \emph{Ramanujan sum} $c_q(a)=\sum_{r\le q,\,(r,q)=1}e(\frac{ar}q)$ equals $\mu(q)$ when $(a,q)=1$: substituting $r\mapsto ar$ we may take $a=1$, and then by $\sum_{d\mid(r,q)}\mu(d)=\mathbf 1_{(r,q)=1}$,
\[
c_q(1)=\sum_{d\mid q}\mu(d)\sum_{s\le q/d}e\Bigl(\frac s{q/d}\Bigr)=\mu(q),
\]
since the inner sum is $0$ unless $\frac qd=1$.

We may take $N$ large. Let $R(t)=\sum_{n\le t}\Lambda(n)e(\frac{an}q)-\frac{\mu(q)}{\varphi(q)}\lfloor t\rfloor$. For $t\le\sqrt N$, $|R(t)|\ll\sqrt N$. For $\sqrt N<t\le N$, the $n$ with $(n,q)>1$ are powers of primes dividing $q$ and contribute at most $\omega(q)\log t\le\mathcal L^2$, with $\mathcal L=\log N$; the others are grouped by their class $r$ modulo $q$:
\[
\sum_{\substack{n\le t\\(n,q)=1}}\Lambda(n)e\Bigl(\frac{an}q\Bigr)=\sum_{(r,q)=1}e\Bigl(\frac{ar}q\Bigr)\psi(t;q,r)=\frac{\mu(q)}{\varphi(q)}t+O\Bigl(q\max_{(r,q)=1}\Bigl|\psi(t;q,r)-\frac t{\varphi(q)}\Bigr|\Bigr).
\]
Since $\log t>\frac12\mathcal L$, Siegel--Walfisz with exponent $A'=A+B+1$ applies for $N$ large ($q\le\mathcal L^B\le(\log t)^{A'}$) and bounds the maximum by $C(A')t(\log t)^{-A'}\le2^{A'}C(A')N\mathcal L^{-A-B-1}$. So $|R(t)|\ll_{A,B}N\mathcal L^{-A}$ for all $t\le N$. Finally, by Abel summation,
\[
\sum_{n\le N}\Bigl(\Lambda(n)e\Bigl(\frac{an}q\Bigr)-\frac{\mu(q)}{\varphi(q)}\Bigr)e(\beta n)=R(N)e(\beta N)-2\pi i\beta\int_1^NR(t)e(\beta t)\dd t,
\]
which is at most $(1+2\pi|\beta|N)\max_{t\le N}|R(t)|$.
\end{proof}

\begin{figure}[ht]
\centering
\fig{ch05-primesum}
\caption{$|S(\alpha)|/\psi(N)$ for $N=10^5$ and $0\le\alpha\le\frac12$ (each bar is the maximum over a window of width $\frac1{4000}$). The peaks sit at the fractions $\frac aq$, with height $\frac{|\mu(q)|}{\varphi(q)}$, as Lemma~\ref{lem:majorarc} predicts: $0.9998$ at $\frac12$, $0.499$ at $\frac13$ and $\frac16$, $0.250$ at $\frac15$, $\frac25$, $\frac1{10}$ and $\frac3{10}$, $0.167$ at $\frac17$, and $0.10$ at $\frac1{11}$. There is no peak at $\frac14$, $\frac18$ or $\frac19$, where $\mu(q)=0$. Between the peaks, Theorem~\ref{thm:vinogradov3} keeps $S(\alpha)$ small.}
\label{fig:primesum}
\end{figure}

\begin{corollary}[Vinogradov, 1937]\label{cor:palpha}
For irrational $\alpha$, the sequence $(p\alpha)$, $p$ running over the primes in increasing order, is equidistributed modulo $1$.
\end{corollary}
\begin{proof}
Fix an integer $j\ne0$; then $\gamma=j\alpha$ is irrational. We show $S(\gamma)=o(N)$. Let $\mathcal L=\log N$ and $Q=N\mathcal L^{-12}$, and choose $\frac aq$ by Lemma~\ref{lem:dirichlet}: $q\le Q$, $|\gamma-\frac aq|\le\frac1{qQ}\le\frac1{q^2}$. If $q>\mathcal L^{12}$, Theorem~\ref{thm:vinogradov3} gives $S(\gamma)\ll(N\mathcal L^{-6}+N^{4/5}+\sqrt{NQ})\mathcal L^5\ll N\mathcal L^{-1}$, since $\sqrt{NQ}=N\mathcal L^{-6}$. If $q\le\mathcal L^{12}$, Lemma~\ref{lem:majorarc} with $B=12$, $A=13$ and $|\beta|N\le\frac N{qQ}\le\mathcal L^{12}$ gives $|S(\gamma)|\le\frac N{\varphi(q)}+O(N\mathcal L^{-1})$.
It remains to see that, in the second case, $q=q_N\to\infty$ as $N\to\infty$. Otherwise, along a sequence of $N$, $q_N\le K$ for a fixed $K$, and the fractions $\frac{a_N}{q_N}$, which tend to $\gamma$, lie in the set of fractions with denominator at most $K$, which has no accumulation point; so $\gamma$ would be one of them, rational. Since $\varphi(q)\ge\sqrt{q/2}$, $\frac N{\varphi(q_N)}=o(N)$.
So $T(t)=\sum_{p\le t}(\log p)e(p\gamma)=S(\gamma)+O(\sqrt t\log^2t)$, with $N=t$, is $o(t)$. By partial summation,
\[
\sum_{p\le N}e(p\gamma)=\frac{T(N)}{\log N}+\int_2^N\frac{T(t)}{t\log^2t}\dd t=o\Bigl(\frac N{\log N}\Bigr),
\]
splitting the integral at $\sqrt N$. As $\pi(N)\asymp\frac N{\log N}$, Weyl's criterion concludes.
\end{proof}

\begin{beauty}
The proof divides the circle into two kinds of points. Near fractions with small denominators, the \emph{major arcs}, the primes are seen through their distribution in progressions, and the sum is large for a reason we can compute. Everywhere else, the \emph{minor arcs}, the bilinear structure of $\Lambda$ forces cancellation. This dichotomy, between structure and randomness, is the circle method of Chapters~6 and~7, and it will return in Chapters~12 and~13 as the dichotomy between sets with a large Fourier coefficient and sets that look random.
\end{beauty}

\section{Vinogradov's mean value theorem}

For polynomial phases $e(\alpha_1n+\cdots+\alpha_kn^k)$, Weyl's differencing loses a square root at each step, so for degree $k$ it saves only about $N^{2^{1-k}}$. Vinogradov's method saves much more, through a count of solutions of a system of equations. For integers $s,k\ge1$ and $N\ge1$, let $J_{s,k}(N)$ be the number of solutions of
\[
x_1^j+\cdots+x_s^j=y_1^j+\cdots+y_s^j\qquad(1\le j\le k)
\]
in integers $1\le x_i,y_i\le N$. With $f(\boldsymbol\alpha)=\sum_{n\le N}e(\alpha_1n+\cdots+\alpha_kn^k)$, orthogonality gives $J_{s,k}(N)=\int_{[0,1]^k}|f(\boldsymbol\alpha)|^{2s}\dd\boldsymbol\alpha$.

\begin{proposition}\label{prop:vmvtbounds}
(a) $J_{s,k}(N)\ge N^s$ and $J_{s,k}(N)\ge s^{-k}N^{2s-k(k+1)/2}$.
(b) If $s\le k$, then $J_{s,k}(N)\le s!\,N^s$.
\end{proposition}
\begin{proof}
(a) The diagonal solutions $y=x$ give $N^s$. For a vector $v\in\Z^k$ let $r(v)$ be the number of $x\in[1,N]^s$ with $(\sum_ix_i^j)_{j\le k}=v$. Then $\sum_vr(v)=N^s$ and $J_{s,k}(N)=\sum_vr(v)^2$. The $j$-th coordinate of $v$ lies in $[1,sN^j]$, so at most $s^kN^{k(k+1)/2}$ vectors $v$ occur, and Cauchy--Schwarz gives $N^{2s}\le s^kN^{k(k+1)/2}J_{s,k}(N)$.
(b) If $s\le k$, a solution has the same power sums $p_j=\sum_ix_i^j$ for $j\le s$ as $y$. By Newton's identities, $p_1,\dots,p_s$ determine the elementary symmetric functions $e_1,\dots,e_s$ of $x_1,\dots,x_s$, hence the polynomial $\prod_i(T-x_i)$, hence the multiset $\{x_i\}$. So $y$ is a permutation of $x$: at most $s!$ choices of $y$ for each $x$.
\end{proof}

The \emph{main conjecture} said that the lower bounds are essentially sharp. It was proved in 2015.

\begin{theorem}[Bourgain, Demeter, Guth; Wooley for $k=3$; quoted]\label{thm:vmvt}
For all $s,k\ge1$ and $\varepsilon>0$, $J_{s,k}(N)\ll_{s,k,\varepsilon}N^{s+\varepsilon}+N^{2s-\frac{k(k+1)}2+\varepsilon}$.
\end{theorem}

Vinogradov himself proved in 1935 a weaker form, sharp for $s$ of size about $k^2\log k$, by an ingenious induction on $p$-adic congruences; Wooley's ``efficient congruencing'' (2012) came close to the conjecture, and Bourgain, Demeter and Guth finished it with \emph{decoupling}, a tool from harmonic analysis about how waves on a curved surface interfere. The theorem gives Weyl sums for degree $k$ a saving of about $N^{1/(k(k-1))}$ instead of $N^{2^{1-k}}$, and its earlier versions gave the Vinogradov--Korobov zero-free region of Chapter~1. We shall not need it: Chapter~7 uses only Hua's lemma, which is proved there by Weyl differencing.

\section{Kloosterman and Salié sums}

For an integer $q\ge1$ and integers $a,b$, the \emph{Kloosterman sum} is
\[
K(a,b;q)=\sum_{\substack{x\bmod q\\(x,q)=1}}e\Bigl(\frac{ax+b\overline x}q\Bigr),
\]
where $\overline x$ is the inverse of $x$ modulo $q$. It counts nothing, but it appears whenever an inverse modulo $q$ appears in a counting problem: in Kloosterman's work on quaternary forms, in the Fourier coefficients of modular forms (the Petersson formula), in the distribution of $\overline x$ itself (Exercise~\ref{ex:inverses}), and in every advance on primes in progressions beyond $\sqrt x$.

\begin{proposition}\label{prop:klbasic}
(a) $K(a,b;q)$ is real. (b) For $(c,q)=1$, $K(ac,b;q)=K(a,bc;q)$; in particular $K(a,b;q)=K(1,ab;q)$ when $(a,q)=1$. (c) $K(a,0;q)=c_q(a)$, the Ramanujan sum; for a prime $p\nmid a$, $K(a,0;p)=-1$.
\end{proposition}
\begin{proof}
(a) The conjugate is $\sum_xe(-\frac{ax+b\overline x}q)$; the substitution $x\mapsto-x$, under which $\overline x\mapsto-\overline x$, turns it back into $K(a,b;q)$. (b) Substitute $y=cx$, so $\overline x=c\overline y$. (c) The first claim is the definition; for $p\nmid a$, $\sum_{x=1}^{p-1}e(\frac{ax}p)=-1$.
\end{proof}

From now on $p$ is an odd prime and we write $K(n)=K(1,n;p)$. A sum of $p-1$ terms of modulus $1$ could be as large as $p-1$. Its mean square is about $p$, so $K(n)$ is typically of size $\sqrt p$, and its fourth moment shows that it is never much larger than $p^{3/4}$.

\begin{lemma}[moments]\label{lem:klmoments}
For an odd prime $p$,
\[
\sum_{n=1}^{p-1}K(n)^2=p^2-p-1,\qquad\sum_{n=1}^{p-1}K(n)^4=2p^3-3p^2-3p-1.
\]
\end{lemma}
\begin{proof}
For the second moment, expand and sum over $n$ first, now including $n=0$:
\[
\sum_{n\bmod p}K(n)^2=\sum_{x,y\ne0}e\Bigl(\frac{x+y}p\Bigr)\sum_ne\Bigl(\frac{n(\overline x+\overline y)}p\Bigr)=p\sum_{x\ne0}e(0)=p(p-1),
\]
since $\overline x+\overline y\equiv0$ exactly when $y=-x$. Subtract $K(0)^2=1$.
For the fourth moment we sum over all pairs $(a,b)$. As $K(a,b;p)$ is real, $K^4=K^2\overline K^2$, and orthogonality in $a$ and in $b$ gives
\[
\sum_{a,b\bmod p}K(a,b;p)^4=p^2\,\#\bigl\{x\in(\mathbb F_p^\times)^4:\ x_1+x_2=x_3+x_4,\ \overline{x_1}+\overline{x_2}=\overline{x_3}+\overline{x_4}\bigr\}.
\]
Count the solutions by $s=x_1+x_2$. If $s=0$, then $x_2=-x_1$, the second equation forces $x_4=-x_3$, and there are $(p-1)^2$ solutions. If $s\ne0$, then $\overline{x_1}+\overline{x_2}=\frac s{x_1x_2}$, so $x_3+x_4=s$ and $x_3x_4=x_1x_2$: $\{x_3,x_4\}$ is the set of roots of $T^2-sT+x_1x_2$, which is $\{x_1,x_2\}$. That gives $2$ ordered choices if $x_1\ne x_2$ and $1$ if $x_1=x_2$. Among the $(p-1)(p-2)$ pairs with $s\ne0$, $p-1$ have $x_1=x_2$ (as $p$ is odd), so this case gives $2(p-1)(p-2)-(p-1)$ solutions. In total $3(p-1)(p-2)$, and $\sum_{a,b}K^4=3p^2(p-1)(p-2)$.
Now sort the pairs. The pair $(0,0)$ gives $K=p-1$; the $2(p-1)$ pairs with exactly one zero give $K=-1$ by Proposition~\ref{prop:klbasic}(c); and for $ab\ne0$, $K(a,b;p)=K(ab)$, each $n\ne0$ arising from $p-1$ pairs. So
\[
(p-1)\sum_{n\ne0}K(n)^4=3p^2(p-1)(p-2)-(p-1)^4-2(p-1),
\]
and dividing by $p-1$ gives $3p^2(p-2)-(p-1)^3-2=2p^3-3p^2-3p-1$.
\end{proof}

\begin{theorem}[Kloosterman, 1926]\label{thm:kloosterman}
For an odd prime $p$ and integers $a,b$ with $p\nmid ab$, $|K(a,b;p)|<2^{1/4}p^{3/4}$.
\end{theorem}
\begin{proof}
$K(a,b;p)=K(ab)$, and one term of a sum of non-negative terms is at most the sum: $K(ab)^4\le2p^3-3p^2-3p-1<2p^3$.
\end{proof}

\begin{theorem}[Weil, 1948; quoted]\label{thm:weil}
For a prime $p$ and $p\nmid ab$, $|K(a,b;p)|\le2\sqrt p$.
\end{theorem}

So one can write $K(n)=2\sqrt p\cos\theta_n$ with $\theta_n\in[0,\pi]$. Weil's proof identifies $e^{\pm i\theta_n}\sqrt p$ as the two eigenvalues of Frobenius acting on the cohomology of a curve, and the bound is the Riemann hypothesis for that curve. Katz showed in 1988 that, as $p\to\infty$, the angles $\theta_n$, $1\le n<p$, become equidistributed with respect to the \emph{Sato--Tate measure} $\frac2\pi\sin^2\theta\dd\theta$; equivalently, $\frac{K(n)}{2\sqrt p}$ follows the semicircle law of Figure~\ref{fig:kloosterman}. Lemma~\ref{lem:klmoments} is consistent with this: the semicircle law has $\mathbb E[(2\cos\theta)^2]=1$ and $\mathbb E[(2\cos\theta)^4]=2$, and the lemma gives $\frac1{p-1}\sum K(n)^2/p\to1$ and $\frac1{p-1}\sum K(n)^4/p^2\to2$.

\begin{figure}[ht]
\centering
\fig{ch05-kloosterman}
\caption{Histogram of the $100002$ normalised Kloosterman sums $\frac{K(1,n;p)}{2\sqrt p}$, $1\le n<p$, for $p=100003$, against the semicircle density $\frac2\pi\sqrt{1-x^2}$. The largest value is $0.99988$, inside Weil's bound $1$; the second and fourth moments of $\frac{K(n)}{\sqrt p}$ are $1.0000$ and $2.0000$, as Lemma~\ref{lem:klmoments} predicts.}
\label{fig:kloosterman}
\end{figure}

When the Legendre symbol is inserted, everything becomes explicit. The \emph{Salié sum} is $S(m,n;p)=\sum_{x\ne0}\bigl(\frac xp\bigr)e(\frac{mx+n\overline x}p)$.

\begin{theorem}[Salié, 1932]\label{thm:salie}
Let $p$ be an odd prime and $p\nmid n$. Then
\[
S(m,n;p)=\Bigl(\frac np\Bigr)g_p\sum_{\substack{y\bmod p\\y^2\equiv mn}}e\Bigl(\frac{2y}p\Bigr).
\]
In particular $|S(m,n;p)|\le2\sqrt p$, and $S(m,n;p)=0$ when $mn$ is a non-residue.
\end{theorem}
\begin{proof}
Let $\Sigma$ be the sum over $y$. Detect $y^2\equiv mn$ by $\frac1p\sum_te(\frac{t(y^2-mn)}p)$:
\[
\Sigma=\frac1p\sum_{t\bmod p}e\Bigl(-\frac{tmn}p\Bigr)\sum_ye\Bigl(\frac{ty^2+2y}p\Bigr).
\]
The term $t=0$ vanishes, since $\sum_ye(\frac{2y}p)=0$. For $t\ne0$, complete the square: $ty^2+2y=t(y+\overline t)^2-\overline t$. And $\sum_ye(\frac{ty^2}p)=\sum_z(1+(\frac zp))e(\frac{tz}p)=(\frac tp)g_p$, because $z$ has $1+(\frac zp)$ square roots. Hence
\[
\Sigma=\frac{g_p}p\sum_{t\ne0}\Bigl(\frac tp\Bigr)e\Bigl(-\frac{tmn+\overline t}p\Bigr).
\]
Substitute $t=-\overline n\,x$: then $-tmn=mx$, $-\overline t=n\overline x$, and $(\frac tp)=(\frac{-1}p)(\frac np)(\frac xp)$. So $\Sigma=\frac{g_p}p(\frac{-1}p)(\frac np)S(m,n;p)$. Multiply by $(\frac np)g_p$ and use $g_p^2=(\frac{-1}p)p$: $(\frac np)g_p\Sigma=S(m,n;p)$. The sum $\Sigma$ has at most two terms of modulus $1$, and $|g_p|=\sqrt p$.
\end{proof}

So the Salié sum is $\pm\sqrt p$ times $0$ or $2\cos\frac{4\pi y}p$, an explicit angle: there is no mystery left, and its angles are equidistributed for the uniform measure, not for Sato--Tate (numerical evaluation shows the difference). For Kloosterman sums no formula of this kind exists, which is why Weil's theorem needs geometry.

\paragraph{Pseudorandom numbers.} A linear congruential generator $x_{n+1}=ax_n+b\pmod m$ is fast, but its successive pairs $(x_n,x_{n+1})$ lie on few parallel lines: Marsaglia's 1968 paper was called ``Random numbers fall mainly in the planes''. The \emph{inversive} generator of Eichenauer-Herrmann and Lehn (1986), $x_{n+1}=a\overline{x_n}+b\pmod p$, has no such lattice structure: its pairs are the points $(x,a\overline x+b)$, and Kloosterman sums measure exactly how evenly they fill the square. Exercise~\ref{ex:inverses} turns Theorem~\ref{thm:kloosterman} into a discrepancy bound, and a numerical calculation runs the two generators through a serial test.

\section{Exercises}

\begin{exercise}[$\star$]\label{ex:ramanujan}
Show that $c_q(a)=\sum_{d\mid(a,q)}d\,\mu(\frac qd)$ for all integers $a$ and $q\ge1$. Deduce that $c_q(a)$ is an integer, and compute $c_{12}(a)$ for $1\le a\le12$.
\end{exercise}

\begin{exercise}[$\star$]\label{ex:sthird}
Show that $\sum_{n\le N}\Lambda(n)e(\frac n3)\sim-\frac N2$ and $\sum_{n\le N}\Lambda(n)e(\frac n4)=o(N)$, directly from the prime number theorem in progressions modulo $3$ and $4$.
\end{exercise}

\begin{exercise}[$\star$]\label{ex:salie}
Let $p$ be an odd prime, $p\nmid n$, and $mn\equiv y_0^2\not\equiv0$. Show that $S(m,n;p)=\pm2\sqrt p\cos\frac{4\pi y_0}p$ when $p\equiv1\pmod4$, and $\pm2i\sqrt p\cos\frac{4\pi y_0}p$ when $p\equiv3\pmod4$, for some sign. What is $S(0,n;p)$?
\end{exercise}

\begin{exercise}[$\star\star$]\label{ex:twisted}
Let $q=q_1q_2$ with $(q_1,q_2)=1$, and let $\overline{q_2}$ be an inverse of $q_2$ modulo $q_1$ and $\overline{q_1}$ one of $q_1$ modulo $q_2$. Show that
\[
K(a,b;q)=K(a\overline{q_2},b\overline{q_2};q_1)\,K(a\overline{q_1},b\overline{q_1};q_2).
\]
Deduce from Weil's bound that $|K(a,b;q)|\le2^{\omega(q)}\sqrt q$ for squarefree $q$ with $(ab,q)=1$.
\end{exercise}

\begin{exercise}[$\star\star$]\label{ex:klpz}
Using Lemma~\ref{lem:klmoments} and the Paley--Zygmund inequality, show that for every large prime $p$, more than a quarter of the $n\in[1,p-1]$ satisfy $|K(1,n;p)|\ge\frac12\sqrt p$. So the bound $\sqrt p$ is the truth for most $n$, not only an upper bound.
\end{exercise}

\begin{exercise}[$\star\star$]\label{ex:inverses}
Let $p$ be an odd prime and $I,J\subseteq[1,p-1]$ intervals of integers. Show that the number of $x\in I$ with $\overline x\bmod p\in J$ is $\frac{|I||J|}p+O(p^{3/4}\log^2p)$, with an absolute constant. What would Weil's bound give?
\end{exercise}

\begin{exercise}[$\star\star\star$]\label{ex:mobiussum}
(Davenport, 1937.) Check the identity $\frac1\zeta=2G-\zeta G^2+(\frac1\zeta-G)(1-\zeta G)$ for $G(s)=\sum_{b\le U}\mu(b)b^{-s}$, and adapt the proof of Theorem~\ref{thm:vinogradov3} to show that, under the same hypotheses,
\[
\sum_{n\le N}\mu(n)e(n\alpha)\ll\Bigl(\frac N{\sqrt q}+N^{5/6}+\sqrt{Nq}\Bigr)(\log N)^5.
\]
\end{exercise}

\section{Hints}

\begin{hint}{ex:ramanujan}
Use $\mathbf 1_{(r,q)=1}=\sum_{d\mid(r,q)}\mu(d)$, and $\sum_{s\le m}e(\frac{as}m)=m\mathbf 1_{m\mid a}$.
\end{hint}

\begin{hint}{ex:sthird}
$e(\frac n3)$ takes the value $-\frac12\pm\frac{\sqrt3}2i$ on the classes $1$ and $2$ modulo $3$.
\end{hint}

\begin{hint}{ex:salie}
$\overline{g_p}=(\frac{-1}p)g_p$ and $g_p^2=(\frac{-1}p)p$, so $g_p=\pm\sqrt p$ or $\pm i\sqrt p$.
\end{hint}

\begin{hint}{ex:twisted}
By the Chinese remainder theorem, write $x=x_1q_2+x_2q_1$ with $x_1$ modulo $q_1$, $x_2$ modulo $q_2$.
\end{hint}

\begin{hint}{ex:klpz}
For $Z\ge0$ and $0\le\theta\le1$, $P(Z>\theta\mathbb EZ)\ge(1-\theta)^2\frac{(\mathbb EZ)^2}{\mathbb EZ^2}$. Take $Z=K(n)^2$ with $n$ uniform.
\end{hint}

\begin{hint}{ex:inverses}
Expand both indicator functions in additive characters modulo $p$; the count becomes $\frac1{p^2}\sum_{a,b}\widehat{\mathbf 1_I}(a)\widehat{\mathbf 1_J}(b)K(a,b;p)$.
\end{hint}

\begin{hint}{ex:mobiussum}
The three pieces are $2\mu(n)\mathbf 1_{n\le U}$, a type I sum with coefficients bounded by $\tau(t)$, $t\le U^2$, and a type II sum with $\mu(k)$, $k>U$, against $c(m)$, $m>U$. Take $U=N^{1/3}$.
\end{hint}

\section{Solutions}

\begin{solution}{ex:ramanujan}
\idea Möbius inversion turns the condition ``$r$ prime to $q$'' into complete geometric sums, which are $0$ or the length.
\writeup By $\sum_{d\mid(r,q)}\mu(d)=\mathbf 1_{(r,q)=1}$,
\[
c_q(a)=\sum_{d\mid q}\mu(d)\sum_{s\le q/d}e\Bigl(\frac{as}{q/d}\Bigr)=\sum_{d\mid q}\mu(d)\frac qd\mathbf 1_{\frac qd\mid a}.
\]
Writing $d'=\frac qd$, this is $\sum_{d'\mid(a,q)}d'\mu(\frac q{d'})$, an integer. For $q=12$, $c_{12}(a)$ depends only on $(a,12)$, and $\mu(12)=\mu(4)=0$, $\mu(6)=1$, $\mu(3)=\mu(2)=-1$. The formula gives $0$ when $(a,12)=1$ or $3$; $2\mu(6)=2$ when $(a,12)=2$; $2\mu(6)+4\mu(3)=-2$ when $(a,12)=4$; $2\mu(6)+6\mu(2)=-4$ when $(a,12)=6$; and $2-4-6+12=4=\varphi(12)$ when $(a,12)=12$. So $c_{12}(a)$ for $a=1,\dots,12$ is $0,2,0,-2,0,-4,0,-2,0,2,0,4$; they sum to $0$, as they must.
\end{solution}

\begin{solution}{ex:sthird}
\idea Group the primes by their class: the exponential becomes a constant on each class, and the prime number theorem in progressions gives the weights.
\writeup Modulo $3$, $e(\frac n3)=1$ for $n\equiv0$, $\omega=-\frac12+\frac{\sqrt3}2i$ for $n\equiv1$ and $\overline\omega$ for $n\equiv2$. The $n\equiv0$ contribute $\sum_{3^k\le N}\log3=O(\log N)$, and $\psi(N;3,1)\sim\psi(N;3,2)\sim\frac N2$, so the sum is $(\omega+\overline\omega)\frac N2+o(N)=-\frac N2+o(N)$. Modulo $4$, $e(\frac n4)$ is $i$ on $n\equiv1$, $-i$ on $n\equiv3$, and the even $n$ contribute $O(\log N)$; since $\psi(N;4,1)\sim\psi(N;4,3)\sim\frac N2$, the sum is $o(N)$. This is Lemma~\ref{lem:majorarc} with $\beta=0$: $\frac{\mu(3)}{\varphi(3)}=-\frac12$ and $\mu(4)=0$.
\end{solution}

\begin{solution}{ex:salie}
\idea Theorem~\ref{thm:salie} leaves only the Gauss sum and two roots of unity, which pair into a cosine.
\writeup By Exercise~3.3 of Tome~5, $\overline{g_p}=(\frac{-1}p)g_p$ and $|g_p|^2=p$, so $g_p$ is real, $\pm\sqrt p$, when $p\equiv1\pmod4$, and purely imaginary, $\pm i\sqrt p$, when $p\equiv3\pmod4$. The roots of $y^2\equiv mn$ are $\pm y_0$, so the sum over $y$ is $e(\frac{2y_0}p)+e(-\frac{2y_0}p)=2\cos\frac{4\pi y_0}p$. With the sign $(\frac np)$, Theorem~\ref{thm:salie} gives the two formulas. For $m=0$, the only root is $y=0$, so $S(0,n;p)=(\frac np)g_p$; directly, $S(0,n;p)=\sum_x(\frac{\overline x}p)e(\frac{n\overline x}p)=(\frac np)g_p$.
\end{solution}

\begin{solution}{ex:twisted}
\idea The Chinese remainder theorem splits $x$, its inverse and the phase into a $q_1$-part and a $q_2$-part.
\writeup As $x_1$ runs over the units modulo $q_1$ and $x_2$ over the units modulo $q_2$, $x=x_1q_2+x_2q_1$ runs over the units modulo $q$. The inverse of $x$ modulo $q$ is $y=y_1q_2+y_2q_1$ with $y_1\equiv\overline{x_1}\,\overline{q_2}^{\,2}\pmod{q_1}$ and $y_2\equiv\overline{x_2}\,\overline{q_1}^{\,2}\pmod{q_2}$: indeed $xy\equiv x_1y_1q_2^2\equiv1\pmod{q_1}$, and similarly modulo $q_2$. Now $e(\frac c{q_1q_2})$ with $c=c_1q_2+c_2q_1$ equals $e(\frac{c_1}{q_1})e(\frac{c_2}{q_2})$. With $c=ax+by$, $c_1=ax_1+by_1$ and $c_2=ax_2+by_2$, so
\[
K(a,b;q)=\sum_{x_1}e\Bigl(\frac{ax_1+b\overline{q_2}^{\,2}\overline{x_1}}{q_1}\Bigr)\sum_{x_2}e\Bigl(\frac{ax_2+b\overline{q_1}^{\,2}\overline{x_2}}{q_2}\Bigr).
\]
By Proposition~\ref{prop:klbasic}(b) with $c=\overline{q_2}$, $K(a,b\overline{q_2}^{\,2};q_1)=K(a\overline{q_2},b\overline{q_2};q_1)$, and the same for $q_2$. For squarefree $q$ with $(ab,q)=1$, induction on the number of prime factors and Weil's bound at each prime give $|K(a,b;q)|\le\prod_{p\mid q}2\sqrt p=2^{\omega(q)}\sqrt q$.
\end{solution}

\begin{solution}{ex:klpz}
\idea Two moments pin down a distribution enough to keep a fixed proportion of it away from $0$.
\writeup Let $n$ be uniform on $[1,p-1]$ and $Z=K(n)^2$. By Lemma~\ref{lem:klmoments}, $\mathbb EZ=\frac{p^2-p-1}{p-1}=p-\frac1{p-1}$ and $\mathbb EZ^2=\frac{2p^3-3p^2-3p-1}{p-1}\le2p^2$. Let $\theta=\frac p{4\mathbb EZ}$, which tends to $\frac14$. The Paley--Zygmund inequality (from Cauchy--Schwarz: $\mathbb E[Z\mathbf 1_{Z>\theta\mathbb EZ}]\ge(1-\theta)\mathbb EZ$ and $\mathbb E[Z\mathbf 1_{Z>\theta\mathbb EZ}]^2\le\mathbb EZ^2P(Z>\theta\mathbb EZ)$) gives
\[
P\Bigl(K(n)^2>\frac p4\Bigr)\ge(1-\theta)^2\frac{(\mathbb EZ)^2}{\mathbb EZ^2}\ge(1-\theta)^2\frac{(p-1)^2}{2p^2}\longrightarrow\frac9{32}>\frac14.
\]
So for large $p$, more than a quarter of the $n$ have $|K(n)|>\frac12\sqrt p$.
\end{solution}

\begin{solution}{ex:inverses}
\idea Counting points of a curve in a box is a double Fourier expansion; the zero frequency gives the area, and every other frequency is a Kloosterman sum.
\writeup For an interval $I$ put
\[
\hat I(a)=\sum_{u\in I}e\!\left(-\frac{au}p\right).
\]
Then $\mathbf 1_I(x)=\frac1p\sum_{a\bmod p}\hat I(a)e(\frac{ax}p)$ for $x\in[0,p-1]$. By the bound for geometric sums, $|\hat I(a)|\le\frac1{2\|a/p\|}$ for $a\not\equiv0$, so $\sum_{a\ne0}|\hat I(a)|\le\sum_{1\le k\le p/2}\frac pk\le p(1+\log p)$. The count is
\[
\sum_{x\ne0}\mathbf 1_I(x)\mathbf 1_J(\overline x)=\frac1{p^2}\sum_{a,b}\hat I(a)\hat J(b)K(a,b;p).
\]
The term $a=b=0$ is $\frac{|I||J|(p-1)}{p^2}=\frac{|I||J|}p+O(1)$. The terms with exactly one of $a,b$ zero have $K=-1$ and total at most $\frac1{p^2}(|I|+|J|)p(1+\log p)\le2(1+\log p)$. The terms with $ab\ne0$ total at most $\frac1{p^2}2^{1/4}p^{3/4}(p(1+\log p))^2=2^{1/4}p^{3/4}(1+\log p)^2$ by Theorem~\ref{thm:kloosterman}. With Weil's bound the error becomes $O(\sqrt p\log^2p)$, which is non-trivial for boxes of area much larger than $p^{3/2}\log^2p$, that is, of side larger than $p^{3/4+\varepsilon}$.
\end{solution}

\begin{solution}{ex:mobiussum}
\idea The same rearrangement works for $\mu$: a short piece, a type I piece with divisor-bounded coefficients, and a type II piece; only the balance of $U$ changes, because the type I piece now reaches $U^2$ with coefficients that need Cauchy--Schwarz.
\writeup \emph{The identity.} $(\frac1\zeta-G)(1-\zeta G)=\frac1\zeta-2G+\zeta G^2$, so the right side is $\frac1\zeta$. Comparing coefficients, $\mu(n)=b_1(n)+b_2(n)+b_3(n)$ with $b_1(n)=2\mu(n)\mathbf 1_{n\le U}$, $b_2(n)=-\sum_{t\mid n,\,t\le U^2}\lambda'_t$ where $\lambda'_t=\sum_{bc=t,\,b,c\le U}\mu(b)\mu(c)$, so $|\lambda'_t|\le\tau(t)$, and $b_3(n)=\sum_{km=n,\,k>U,\,m>U}\mu(k)c(m)$ with $c(m)=\sum_{b\mid m,\,b>U}\mu(b)$, as in Lemma~\ref{lem:vaughan}: the coefficient of $m^{-s}$ in $1-\zeta G$ is $c(m)$, which vanishes for $m\le U$, and $\frac1\zeta-G=\sum_{k>U}\mu(k)k^{-s}$.
\emph{The pieces.} Let $\mathcal L=\log N$ and $q\le N$. First, $|\sum_{n\le N}b_1(n)e(n\alpha)|\le2U$. Next, $\sum_{n\le N}b_2(n)e(n\alpha)=-\sum_{t\le U^2}\lambda'_t\sum_{d\le N/t}e(\alpha td)$ is at most $\sum_{t\le U^2}\tau(t)\min(\frac Nt,\frac1{\|\alpha t\|})$. By Cauchy--Schwarz, bounding the minimum by $\frac Nt$ in one factor,
\[
\sum_{t\le U^2}\tau(t)\min\Bigl(\frac Nt,\frac1{\|\alpha t\|}\Bigr)\le\Bigl(N\sum_{t\le U^2}\frac{\tau(t)^2}t\Bigr)^{\frac12}\Bigl(\sum_{t\le U^2}\min\Bigl(\frac Nt,\frac1{\|\alpha t\|}\Bigr)\Bigr)^{\frac12}.
\]
As in Lemma~\ref{lem:phisum}, $\sum_{t\le X}\frac{\tau(t)^2}t\le\sum_{t\le X}\frac{\tau_4(t)}t\le(1+\log X)^4$, and Lemma~\ref{lem:summin2} bounds the second factor. So this piece is $\ll(N\mathcal L^4)^{1/2}((\frac Nq+U^2+q)\mathcal L)^{1/2}\ll(\frac N{\sqrt q}+\sqrt NU+\sqrt{Nq})\mathcal L^{5/2}$. Finally, the type II argument of Theorem~\ref{thm:vinogradov3} applies word for word to $b_3$, with $\mu(k)$ in place of $\Lambda(k)$ (so $\mathcal L^2$ is replaced by $1$) and $V=U$: blocks $m\in\left]M,2M\right]$ with $U\le M<\frac NU$ give $\ll(\frac N{\sqrt q}+\frac N{\sqrt U}+\sqrt{Nq})\mathcal L^5$.
\emph{The choice.} The total is $\ll(\frac N{\sqrt q}+U+\sqrt NU+\frac N{\sqrt U}+\sqrt{Nq})\mathcal L^5$, and $U=N^{1/3}$ makes $\sqrt NU=\frac N{\sqrt U}=N^{5/6}$. Davenport combined this with a major arc analysis, as in Lemma~\ref{lem:majorarc} but with Siegel--Walfisz for $\mu$, to show that $\sum_{n\le N}\mu(n)e(n\alpha)\ll_AN(\log N)^{-A}$ uniformly in $\alpha$.
\end{solution}

\section*{Sources and further reading}

\begin{itemize}
\item \emph{Vinogradov's method.} I.~M.~Vinogradov, ``Representation of an odd number as a sum of three primes'', \emph{C.\ R.\ (Doklady) Acad.\ Sci.\ URSS} 15 (1937), 291--294; \emph{The Method of Trigonometrical Sums in the Theory of Numbers}, translated by K.~F.~Roth and A.~Davenport, Interscience, 1954. R.~C.~Vaughan, ``Sommes trigonométriques sur les nombres premiers'', \emph{C.\ R.\ Acad.\ Sci.\ Paris Sér.\ A} 285 (1977), 981--983. Our proof follows H.~Davenport, \emph{Multiplicative Number Theory}, 3rd ed., Springer GTM~74, 2000, Chapters~24--25, and R.~C.~Vaughan, \emph{The Hardy--Littlewood Method}, 2nd ed., Cambridge, 1997, Chapter~3.
\item \emph{Möbius.} H.~Davenport, ``On some infinite series involving arithmetical functions~II'', \emph{Quart.\ J.\ Math.\ Oxford} 8 (1937), 313--320.
\item \emph{Mean value theorem.} I.~M.~Vinogradov, ``New estimates for Weyl sums'', \emph{Dokl.\ Akad.\ Nauk SSSR} 8 (1935), 195--198. T.~D.~Wooley, ``Vinogradov's mean value theorem via efficient congruencing'', \emph{Ann.\ of Math.}\ 175 (2012), 1575--1627; ``The cubic case of the main conjecture in Vinogradov's mean value theorem'', \emph{Adv.\ Math.}\ 294 (2016), 532--561. J.~Bourgain, C.~Demeter and L.~Guth, ``Proof of the main conjecture in Vinogradov's mean value theorem for degrees higher than three'', \emph{Ann.\ of Math.}\ 184 (2016), 633--682.
\item \emph{Kloosterman and Salié.} H.~D.~Kloosterman, ``On the representation of numbers in the form $ax^2+by^2+cz^2+dt^2$'', \emph{Acta Math.}\ 49 (1926), 407--464. H.~Salié, ``Über die Kloostermanschen Summen $S(u,v;q)$'', \emph{Math.\ Z.}\ 34 (1932), 91--109. A.~Weil, ``On some exponential sums'', \emph{Proc.\ Nat.\ Acad.\ Sci.\ USA} 34 (1948), 204--207. N.~M.~Katz, \emph{Gauss Sums, Kloosterman Sums, and Monodromy Groups}, Annals of Mathematics Studies 116, Princeton, 1988. H.~Iwaniec and E.~Kowalski, \emph{Analytic Number Theory}, AMS Colloquium Publications 53, 2004, Chapter~11.
\item \emph{Pseudorandom generators.} G.~Marsaglia, ``Random numbers fall mainly in the planes'', \emph{Proc.\ Nat.\ Acad.\ Sci.\ USA} 61 (1968), 25--28. J.~Eichenauer and J.~Lehn, ``A non-linear congruential pseudo random number generator'', \emph{Statist.\ Hefte} 27 (1986), 315--326.
\end{itemize}
